\section{项目经历}

\entry{基于 Kafka 消息队列的订单系统异步化重构}{研发工程师}{2025.01 -- 2025.02}{}
\begin{itemize}
  \item \textbf{项目背景：}订单处理系统在高并发场景下存在明显性能瓶颈，影响用户体验与业务增长；负责设计并落地订单异步处理机制，提升系统性能和可扩展性。
  \item \textbf{异步解耦：}引入 Kafka 作为消息队列，将订单流程中的非核心同步操作改造为异步处理，解耦订单服务与其他微服务之间的调用依赖，缩短核心链路并提升系统响应速度和并发处理能力。
  \item \textbf{事件驱动：}采用事件驱动架构，将订单状态变更封装为事件并通知相关服务，降低服务间耦合度，保障订单数据的一致性与处理可靠性。
  \item \textbf{项目成果：}订单处理时间减少 40\%，并发处理能力提升 60\%；部分场景下订单流程耗时缩短至原来的 1\%。
\end{itemize}

\vspace{1pt}
\entry{基于 Redis 的令牌桶限流 SDK}{研发工程师}{2024.03 -- 2024.04}{}
\begin{itemize}
  \item \textbf{项目背景：}高并发场景下的请求洪峰易造成服务不可用和用户体验下降；负责设计通用限流策略，保障系统稳定及请求的公平性与可用性。
  \item \textbf{限流实现：}采用令牌桶算法控制订单请求速率，并使用 Redis 存储共享限流状态，实现分布式限流及限流数据的高可用与一致性。
  \item \textbf{熔断保护：}系统过载时快速失败，阻断过量请求持续冲击下游服务，避免系统过载和故障扩散。
  \item \textbf{项目成果：}请求成功率提升 30\%，用户投诉减少 50\%；核心业务系统未再因流量洪峰出现服务重启或宕机，服务可用性得到提升。
\end{itemize}
